\documentclass[12pt]{article}
\usepackage[utf8]{inputenc}
\usepackage{amsmath}
\usepackage{amsfonts}
\usepackage{amssymb}
\usepackage{graphicx}
\usepackage[left=3.0cm,right=3.0cm,top=3.0cm,bottom=3.0cm]{geometry}
\usepackage[numbers]{natbib}
\usepackage{float}
\usepackage{xcolor}
\usepackage{enumitem}
\usepackage{hyperref}
\setlength{\parskip}{1em}

% ============================================================
%  Exposition — Lecture 04, CCBS (EST 5053), Fall 2026
%  Álvaro Bátrez
%
%  AUTHOR'S NOTE:
%  All figure slots are marked with
%      % === FIGURE n: description ===
%  and use provisional file names inside figs/.
%  Replace them with the graphics/animations exported from
%  Mathematica (PDF or PNG format).
% ============================================================

\title{Lecture 04 --- Dynamics: local stability, phase portraits and bifurcations\\
\large From biochemical structure to dynamical function}
\author{Álvaro Bátrez}
\date{September 10, 2026}

\begin{document}

\maketitle

\begin{abstract}
\noindent A simulation tells us what \emph{one} cell did \emph{once}. In this session
we build the map that goes from the \textbf{structure} of a biochemical network to
its dynamical \textbf{function}: the dimension of the system decides which behaviors
are possible, and chemistry decides which vector fields a cell can have. We walk
through the phase line in one dimension, the construction of memory through three
crossings, the saddle-node bifurcation with its experimental signatures (hysteresis
and \emph{critical slowing down}), the genetic toggle switch in two dimensions,
linearization and the Hurwitz rule, Lyapunov functions for the cases where the
Jacobian stays silent, and the restrictions that chemistry imposes on the
mathematically possible vector fields.
\end{abstract}

\tableofcontents

\newpage

% ============================================================
\section{Why analyze, if we can already simulate?}
% ============================================================

\subsection{Starting point: a model we already master}

In the previous lecture we wrote reactions, assembled the system
$\dot{x} = S\,v(x)$ and integrated the equations. Take the simplest example:
a constitutively produced protein with constant production rate $\mu$, removed
(degradation + growth dilution) at the first-order rate $\lambda$:

\begin{equation}
\dot{x} = \mu - \lambda x,
\qquad
x^{*} = \frac{\mu}{\lambda},
\qquad
x(t) = x^{*} + \left(x_0 - x^{*}\right) e^{-\lambda t}.
\label{eq:constitutive}
\end{equation}

Everything is solved exactly: the steady level is $x^{*} = \mu/\lambda$ and the
relaxation time is $1/\lambda$. Moving a slider in the simulation changes
\emph{both things at once}.

% === FIGURE 1: simulation of equation (1) for 2-3 values of lambda ===
% Mathematica hint: Plot of x(t) with lambda = 0.25, 0.5, 1.0, mu = 2, x0 = 0.4
\begin{figure}[H]
\centering
\includegraphics[width=0.75\textwidth]{figs/fig01_relaxation.pdf}
\caption{Exact solution of equation~\eqref{eq:constitutive} for several values
of $\lambda$. The steady state moves as $1/\lambda$, and so does the relaxation
time.}
\label{fig:relaxation}
\end{figure}

\textbf{Room check.} If $\lambda$ goes from $0.5$ to $1.0$, where does the steady
level move, and by what factor does the relaxation time change? Can you name one
biological change that the slider does \emph{not} represent?

\subsection{Limit 1: the discarded mechanism is not in the run}

Every cellular model is built by \emph{compression}: several promoter binding states
become one effective Hill curve; transcription and translation become a single
production arrow; a conserved pool may drop out of the coordinates. Each compression
says ``this hidden process is fast, fixed, or irrelevant.'' A simulation cannot tell
us what the discarded process would have changed, because that process \emph{is not
in the equations being run}.

\subsection{Limit 2: the family is too large to enumerate}

With ten uncertain parameters sampled at only ten settings each, we already face
$10^{10}$ combinations, before sampling a single initial condition. We need
statements like ``every field with this sign pattern is stable'' or ``no scalar
autonomous system can oscillate'': each one covers \emph{infinitely many} numerical
runs at once.

\subsection{The map we want: from structure to function}

\begin{center}
\begin{tabular}{lll}
\textbf{Structure} (coarse to fine) & $\longrightarrow$ & \textbf{Function} (coarse to fine) \\
\hline
interaction graph & & admissible region; bounded? \\
signs and invariant pools & & stable fixed point? \\
response-curve shapes & & how many states? switch or cycle? \\
parameters + initial state & & period, amplitude, one trajectory \\
\end{tabular}
\end{center}

The bottom arrow is ordinary simulation. The upper arrows are the ones we want to
justify, because each holds for a \emph{family} of models.

% === FIGURE 2: diagram of the structure -> function map (TikZ or Mathematica) ===
\begin{figure}[H]
\centering
\includegraphics[width=0.85\textwidth]{figs/fig02_map.pdf}
\caption{The structure-to-function map that organizes the whole lecture.}
\label{fig:map}
\end{figure}

\begin{quote}
\textbf{Structural does not mean vague.} A structural claim names the information
it uses and the family over which the conclusion holds.
\end{quote}

\newpage

% ============================================================
\section{One dimension: a fixed point is a crossing}
% ============================================================

\subsection{Addition against removal}

Let $x \geq 0$ be a concentration. Write its net rate as an \emph{addition} rate
minus a \emph{removal} rate, both nonnegative:

\begin{equation}
\dot{x} = f^{+}(x) - f^{-}(x) = f(x),
\qquad f^{+}(x) \geq 0,\quad f^{-}(x) \geq 0.
\label{eq:addrem}
\end{equation}

A fixed point $x^{*}$ is a \textbf{crossing} of the addition and removal curves:
$f^{+}(x^{*}) = f^{-}(x^{*})$. Just to the right of the crossing, removal must win
for motion to point back; just to the left, addition must win. At a simple crossing
this is the slope condition

\begin{equation}
f'(x^{*}) = (f^{+})'(x^{*}) - (f^{-})'(x^{*}) < 0.
\label{eq:slope}
\end{equation}

This split uses the physical domain of a concentration ($x \geq 0$, both terms
nonnegative). It is not cosmetic: it is a structural restriction we will return to
in section~\ref{sec:chemistry}.

\subsection{The phase-line recipe}

\begin{enumerate}[nosep]
  \item Plot $f(x)$ (or plot $f^{+}$ and $f^{-}$ separately and compare them).
  \item Mark every zero of $f$.
  \item On each interval between zeros: arrow right if $f>0$, left if $f<0$.
  \item Translate the phase line into time traces $x(t)$ from several initial conditions.
\end{enumerate}

\subsection{Three fields, three behaviors}

\paragraph{(a) Decay: $\dot{x} = -x$.} The field crosses zero with negative slope;
both arrows point to the origin; every trajectory decays exponentially.

\paragraph{(b) Nonhyperbolic growth: $\dot{x} = x^{2}$.} The field only \emph{touches}
zero: it attracts from the left and repels to the right, and $f'(0)=0$ cannot
classify it. Separating variables:

\begin{equation}
x(t) = \frac{x_0}{1 - x_0 t},
\label{eq:blowup}
\end{equation}

so for $x_0 > 0$ there is finite-time escape, at $t = 1/x_0$. The phase line got the
\emph{direction} right, but it carried no clock.

\paragraph{(c) Bistable switch: $\dot{x} = -(x-1)(x-2)(x-3)$.} Two attracting roots
separated by one repelling threshold: memory without a cycle.

% === FIGURE 3: three rows: f(x), phase line, x(t) for the three fields ===
% Mathematica hint: 3x3 Grid with Plot, phase-line arrows, and NDSolve solutions
% for several x0.
\begin{figure}[H]
\centering
\includegraphics[width=0.95\textwidth]{figs/fig03_three_fields.pdf}
\caption{Three coordinated views of each scalar field: $f(x)$, the phase line, and
trajectories $x(t)$. Filled: attractors; open: repellors.}
\label{fig:threefields}
\end{figure}

\subsection{The first impossibility theorem}

Between two consecutive roots the sign of $f$ cannot change, so a nonstationary
scalar trajectory is \textbf{monotone} there. To return to an earlier value it would
have to reverse direction, which requires crossing a zero; but a unique trajectory
cannot cross a fixed point. Therefore:

\begin{quote}
\centering
\emph{A scalar autonomous system can switch, but it cannot sustain a nonconstant
cycle.}
\end{quote}

This is our first structural result: deform the curves without changing the number,
order and transverse character of their crossings, and the same phase portrait
survives. The exact trajectories change; the qualitative basins do not.

\newpage

% ============================================================
\section{To remember, cross three times}
% ============================================================

\subsection{Memory as a design problem}

Memory is a \emph{construction problem}. Two persistent states require the crossing
sequence \textbf{stable -- unstable -- stable}. Consider a positively autoregulated
gene: its own protein activates its production. Compressing the fast redistribution
among promoter states into an activating Hill fraction, with first-order removal
(dilution), a dimensionless teaching model is

\begin{equation}
\frac{dx}{d\tau} =
\underbrace{0.15 + \frac{3x^{4}}{1+x^{4}}}_{f^{+}(x):\ \text{basal + activated production}}
\;-\;
\underbrace{x}_{f^{-}(x):\ \text{dilution removal}}.
\label{eq:autoreg}
\end{equation}

The Hill term is a \emph{lumped} effective law: it assumes promoter dynamics are fast
relative to protein accumulation; the exponent represents a cooperative occupancy
model, \emph{not} an elementary four-molecule collision.

\subsection{The logic of shape}

Design does not start with the numbers; it starts with the \emph{shape}:

\begin{enumerate}[nosep]
  \item \textbf{Positive feedback}: production rises with $x$.
  \item \textbf{Enough steepness} (effective cooperativity): production locally
        outruns removal.
  \item \textbf{Saturation}: production flattens so removal can win again.
\end{enumerate}

The calculated crossings lie near $0.15$, $0.68$ and $3.12$: the outer two attract
and the middle one is the \textbf{threshold}. A transient input is \emph{remembered}
only if it pushes the state across that threshold. In this construction,
\textbf{memory is bought with a sufficiently steep cooperative response}.

% === FIGURE 4: f+ and f- curves of model (4) with the three crossings marked ===
\begin{figure}[H]
\centering
\includegraphics[width=0.8\textwidth]{figs/fig04_three_crossings.pdf}
\caption{Binary memory is three addition--removal crossings. Parameters place the
crossings; shape creates the possibility.}
\label{fig:threecrossings}
\end{figure}

\textbf{Design check.} Sketch a monotonically increasing addition curve that never
becomes steeper than the removal line: can it cross three times? Now make it steep
but remove saturation: which crossing is lost?

\newpage

% ============================================================
\section{When two curves touch, a switch is born}
% ============================================================

\subsection{The saddle-node bifurcation}

Let $r$ be a control parameter and $x$ a local displacement from the collision point.
The \textbf{normal form} of the saddle-node bifurcation is

\begin{equation}
\dot{x} = f(x;r) = r - x^{2}.
\label{eq:saddlenode}
\end{equation}

In the addition--removal language, the two curves \emph{touch} when both their values
\emph{and} their slopes agree:

\begin{equation}
f^{+}(x^{*};r^{*}) = f^{-}(x^{*};r^{*}),
\qquad
\frac{\partial f^{+}}{\partial x}(x^{*};r^{*}) = \frac{\partial f^{-}}{\partial x}(x^{*};r^{*}).
\label{eq:tangency}
\end{equation}

For $r>0$ the roots are $x^{*} = \pm\sqrt{r}$; since $f_x = -2x$, the positive one is
stable and the negative one unstable. At $r=0$ they collide; for $r<0$ neither
exists. A bifurcation is this change in the \emph{organization of all trajectories},
not merely a bend in one plotted path.

% === FIGURE 5: f(x;r) for r<0, r=0, r>0 (three panels) ===
\begin{figure}[H]
\centering
\includegraphics[width=0.95\textwidth]{figs/fig05_saddlenode_fields.pdf}
\caption{The field $f(x;r)=r-x^2$ below, at, and above tangency.}
\label{fig:snfields}
\end{figure}

% === FIGURE 6: bifurcation diagram x* against r (solid and dashed branches) ===
\begin{figure}[H]
\centering
\includegraphics[width=0.7\textwidth]{figs/fig06_saddlenode_diagram.pdf}
\caption{Bifurcation diagram: stable branch $x^{*}=+\sqrt{r}$ (solid) and unstable
branch $x^{*}=-\sqrt{r}$ (dashed). No equilibrium exists for $r<0$.}
\label{fig:sndiagram}
\end{figure}

\subsection{Two experimental signatures}

\paragraph{Hysteresis.} A full sigmoidal switch usually has \emph{two} folds.
Increasing the input destroys one state at one fold; decreasing it destroys the other
state at the \emph{other} fold. The switching threshold therefore depends on
direction and history.

\paragraph{Critical slowing down.} Near the stable root, with
$u = x - \sqrt{r}$,

\begin{equation}
\dot{u} = -2\sqrt{r}\,u
\quad\Longrightarrow\quad
u(t) \approx u(0)\,e^{-2\sqrt{r}\,t},
\qquad
\tau_{\text{relax}} = \frac{1}{2\sqrt{r}}.
\label{eq:slowing}
\end{equation}

Moving from $r=1$ to $r=0.01$ changes $\tau_{\text{relax}}$ from $0.5$ to $5$: a
hundredfold approach to the fold gives a tenfold slowing. This is a \emph{testable}
prediction (early-warning signals for critical transitions~\cite{scheffer2009}),
provided noise, sampling and the valid parameter range are checked.

% === FIGURE 7: tau_relax against r (log scale) or relaxations for several r ===
\begin{figure}[H]
\centering
\includegraphics[width=0.7\textwidth]{figs/fig07_slowing.pdf}
\caption{The relaxation time diverges like $1/(2\sqrt{r})$ as the fold is approached.}
\label{fig:slowing}
\end{figure}

\newpage

% ============================================================
\section{Two dimensions: a second variable buys going around}
% ============================================================

\subsection{Nullclines: species-by-species balances}

In the plane, trajectories can \emph{go around} fixed points, and a saddle can
separate basins. The \textbf{nullcline} of species $i$ is the set where its total
addition equals its total removal, that is, where its own derivative is zero. Only
an \emph{intersection of all nullclines} is a fixed point.

\subsection{The biological toggle, before the equations}

Two genes, $G_X$ and $G_Y$, produce the repressors $X$ and $Y$; each represses the
other:

\begin{equation}
G_X \rightsquigarrow G_X + X, \quad X \to \varnothing,
\qquad
G_Y \rightsquigarrow G_Y + Y, \quad Y \to \varnothing,
\qquad
Y \dashv G_X, \quad X \dashv G_Y.
\label{eq:togglenet}
\end{equation}

For the arm producing $X$: with fast promoter redistribution (rapid equilibrium),
conserved pool $q_{G_X} = G_X + C_{G_XY}$ and occupancy
$C_{G_XY}/G_X = (Y/K_Y)^{n_Y}$, the free-promoter fraction is

\begin{equation}
\frac{G_X}{q_{G_X}} = \frac{1}{1+(Y/K_Y)^{n_Y}}.
\label{eq:occupancy}
\end{equation}

If composite production is proportional to free promoter, the slow concentrations obey

\begin{equation}
\dot{X} = \frac{\alpha_X}{1+(Y/K_Y)^{n_Y}} - \delta_X X,
\qquad
\dot{Y} = \frac{\alpha_Y}{1+(X/K_X)^{n_X}} - \delta_Y Y.
\label{eq:toggle}
\end{equation}

This is \emph{not} elementary mass action: it is a reduced law for a specific
biological scenario, and its assumptions stay attached.

\subsection{Phase portrait of the symmetric case}

Scaling concentrations and time, the symmetric teaching case is

\begin{equation}
\dot{x} = \frac{4}{1+y^{2}} - x,
\qquad
\dot{y} = \frac{4}{1+x^{2}} - y,
\qquad x,y \geq 0.
\label{eq:togglesym}
\end{equation}

Procedure: (1) $\dot{x}=0$ gives the nullcline $x = 4/(1+y^2)$; (2) $\dot{y}=0$ gives
$y = 4/(1+x^2)$; (3) at every intersection both components vanish; (4) testing one
point per region orients the arrows.

\begin{center}
\begin{tabular}{cccc}
\hline
fixed point & eigenvalues & type & interpretation \\
\hline
$(0.27,\ 3.73)$ & $-0.5,\ -1.5$ & stable node & low $X$ / high $Y$ \\
$(1.38,\ 1.38)$ & $+0.31,\ -2.31$ & saddle & threshold state \\
$(3.73,\ 0.27)$ & $-0.5,\ -1.5$ & stable node & high $X$ / low $Y$ \\
\hline
\end{tabular}
\end{center}

The \textbf{stable manifold of the saddle} (the separatrix) is the basin boundary: a
perturbation switches the cell only if it crosses it. A steep response by itself is
\emph{not yet} bistability. This switch was actually built in
\emph{E.~coli}~\cite{gardner2000}; what transfers is the dynamical claim: two
attractors, their basins, and a threshold that makes history matter.

% === FIGURE 8: nullclines + phase portrait of the toggle with separatrix ===
% Mathematica hint: ContourPlot of the nullclines + StreamPlot of the field
% + separatrix curve integrated backwards from the saddle.
\begin{figure}[H]
\centering
\includegraphics[width=0.85\textwidth]{figs/fig08_toggle.pdf}
\caption{Nullclines and phase portrait of the symmetric toggle~\eqref{eq:togglesym}.
Filled points: expression states; open point: saddle; dashed curve: separatrix.}
\label{fig:toggle}
\end{figure}

% === ANIMATION (optional): trajectories from several (x0,y0) converging ===
% Mathematica hint: Animate/Manipulate with NDSolve + StreamPlot.

\newpage

% ============================================================
\section{The dimension boundary}
% ============================================================

Trajectories cannot cross, but the \emph{force} of that rule depends on dimension:

\begin{itemize}[nosep]
  \item \textbf{One dimension:} motion is ordered between fixed points (monotonicity).
  \item \textbf{Two dimensions:} an isolated closed orbit separates the plane into
        inside and outside. The \textbf{Poincar\'e--Bendixson} theorem turns ``no
        fixed point remains in the limiting set'' into the existence of a periodic
        orbit. The qualitative catalogue is short: nodes, saddles, spirals, centers
        and isolated cycles.
  \item \textbf{Three dimensions:} a curve has room to wind, braid and pass around
        itself without violating uniqueness. Chaos becomes possible.
\end{itemize}

Biological example: the \textbf{repressilator}, a cyclic ring of three repressors,

\begin{equation}
\dot{p}_i = \frac{\alpha}{1+p_{i-1}^{\,n}} - p_i,
\qquad i = 1,2,3 \pmod 3,
\label{eq:repressilator}
\end{equation}

realized synthetically in \emph{E.~coli}~\cite{elowitz2000} and later improved toward
persistent, synchronized oscillations~\cite{potvin2016}. Moral:
\textbf{lack of structure is not a structure} --- permission by geometry does not
explain a mechanism. That is why we need the two structural retreats of the next
sections.

% === FIGURE 9: repressilator oscillations p_i(t) ===
\begin{figure}[H]
\centering
\includegraphics[width=0.8\textwidth]{figs/fig09_repressilator.pdf}
\caption{Sustained oscillations of the repressilator~\eqref{eq:repressilator}
(example: large $\alpha$, $n=2$).}
\label{fig:repressilator}
\end{figure}

\newpage

% ============================================================
\section{Close to a fixed point, every network is linear}
% ============================================================

\subsection{Linearization is a new dynamical system}

Let $\dot{x} = F(x)$ with fixed point $x^{*}$ and displacement $u = x - x^{*}$.
Taylor expansion:

\begin{equation}
\dot{u} = A u + R(u),
\qquad
A = J(x^{*}) = \left.\frac{\partial F}{\partial x}\right|_{x^{*}},
\qquad
\frac{\lVert R(u)\rVert}{\lVert u\rVert} \to 0 \ \ (u\to 0).
\label{eq:taylor}
\end{equation}

The canonical linearization discards the remainder and uses a new displacement
vector $z$:

\begin{equation}
\boxed{\ \frac{d}{dt}\begin{bmatrix} z_1 \\ z_2 \end{bmatrix}
= A \begin{bmatrix} z_1 \\ z_2 \end{bmatrix}\ },
\qquad
A = \begin{bmatrix} a_{11} & a_{12} \\ a_{21} & a_{22} \end{bmatrix}.
\label{eq:linear}
\end{equation}

If $A w = \lambda w$, motion along the eigenvector $w$ has amplitude
$\dot{a} = \lambda a$, and every solution is a combination of modes:

\begin{equation}
z(t) = \sum_j c_j\, e^{\lambda_j t} w_j.
\label{eq:modes}
\end{equation}

The real part of $\lambda_j$ sets the exponential envelope; the imaginary part sets
the rotation. For the fixed point to recover from \emph{every} small perturbation,
\emph{all} modes must decay.

\subsection{The Hurwitz rule and the trace--determinant plane}

\begin{quote}
\textbf{Hurwitz rule.} A real matrix $A$ is Hurwitz if and only if
$\operatorname{Re}\lambda_j(A) < 0$ for every eigenvalue. Then $z=0$ is globally
exponentially stable for~\eqref{eq:linear}, and $x^{*}$ is locally exponentially
stable in the nonlinear system. A single eigenvalue with positive real part proves
instability. A zero real part leaves the nonlinear case undecided.
\end{quote}

For a $2\times2$ matrix, with $\tau = \operatorname{tr}A$, $\Delta = \det A$ and
discriminant $D = \tau^2 - 4\Delta$:

\begin{equation}
\boxed{\ A \text{ is Hurwitz} \iff \tau < 0 \ \text{and}\ \Delta > 0\ }
\label{eq:hurwitz2d}
\end{equation}

Intuition: $\Delta < 0$ means two real eigenvalues of opposite sign (one direction
grows: a saddle); $\Delta > 0$ puts them on the same side or in a conjugate pair,
and $\tau < 0$ then places their real parts on the decaying side.

% === FIGURE 10: (tau, Delta) plane with regions: saddle, nodes, spirals, parabola ===
\begin{figure}[H]
\centering
\includegraphics[width=0.8\textwidth]{figs/fig10_trace_det.pdf}
\caption{The trace--determinant plane. The Hurwitz region is $\tau<0$, $\Delta>0$;
the parabola $\tau^2 = 4\Delta$ separates real from complex eigenvalues; the axis
$\Delta>0$, $\tau=0$ holds the linear centers (the undecided case).}
\label{fig:tracedet}
\end{figure}

\subsection{The Jacobian as a wiring diagram}

For the toggle~\eqref{eq:togglesym}:

\begin{equation}
J(x,y) =
\begin{bmatrix}
-1 & -\dfrac{8y}{(1+y^{2})^{2}} \\[8pt]
-\dfrac{8x}{(1+x^{2})^{2}} & -1
\end{bmatrix}.
\label{eq:jactoggle}
\end{equation}

The negative diagonal is each protein's removal (growth dilution contributes
$-\lambda_i$ to the diagonal); the negative off-diagonal entries are the mutual
repression. In any biochemical model, $A_{ij}$ asks how species $j$ changes the net
rate of species $i$: \textbf{the sign pattern of the Jacobian is the local wiring
diagram}.

\textbf{Scope of the Jacobian.} A Hurwitz Jacobian proves \emph{local} attraction
near \emph{one} hyperbolic fixed point. It does not locate every fixed point, does
not draw the global basin boundary, does not prove boundedness and does not rule out
a remote attractor. At the toggle saddle, an eigenvector only gives the
\emph{tangent} to the separatrix; the global curve still belongs to the nonlinear
system.

% === FIGURE 11: gallery of 4 portraits: node, saddle, spiral, center ===
\begin{figure}[H]
\centering
\includegraphics[width=0.95\textwidth]{figs/fig11_gallery.pdf}
\caption{The four modal portraits: stable node, saddle, stable spiral and linear
center. They are conclusions from eigenvalues, not decorations.}
\label{fig:gallery}
\end{figure}

\newpage

% ============================================================
\section{When the linear test says nothing}
% ============================================================

\subsection{Three systems, one and the same Jacobian}

Compare three systems in polar coordinates $(r,\theta)$:

\begin{equation}
\dot{r} = -r^{3},\ \dot{\theta}=1
\qquad
\dot{r} = 0,\ \dot{\theta}=1
\qquad
\dot{r} = +r^{3},\ \dot{\theta}=1.
\label{eq:centers}
\end{equation}

At the origin, \emph{all three} have Jacobian
$\left(\begin{smallmatrix} 0 & -1 \\ 1 & 0 \end{smallmatrix}\right)$ with eigenvalues
$\pm i$. The cubic radial term is invisible at first order, and it decides three
different fates: attraction, neutral cycles or repulsion.

% === FIGURE 12: the three fields of (19) side by side ===
\begin{figure}[H]
\centering
\includegraphics[width=0.95\textwidth]{figs/fig12_centers.pdf}
\caption{Same Jacobian, three fates. Only the cubic term decides.}
\label{fig:centers}
\end{figure}

\subsection{The spring: state, energy and LaSalle}

The damped spring $m\ddot{q} + c\dot{q} + kq = 0$ has two jobs here.
First, \emph{what is a state?}: position alone is not enough (the mass can pass the
same $q$ moving left or right); with $v=\dot{q}$ and $z=(q,v)$ we get an
\emph{exact} first-order system (not a linearization):

\begin{equation}
\frac{d}{dt}\begin{bmatrix} q \\ v \end{bmatrix}
=
\begin{bmatrix} 0 & 1 \\ -k/m & -c/m \end{bmatrix}
\begin{bmatrix} q \\ v \end{bmatrix}.
\label{eq:spring}
\end{equation}

Second, mechanical energy as a global certificate:

\begin{equation}
E(q,v) = \tfrac{1}{2}mv^{2} + \tfrac{1}{2}kq^{2}
\quad\Longrightarrow\quad
\dot{E} = mv\dot{v} + kq\dot{q} = -c\,v^{2} \leq 0.
\label{eq:energy}
\end{equation}

With $c=0$ energy is conserved: each initial condition lives on its own ellipse; the
origin is stable but \emph{not} attracting. With $c>0$ energy never increases.
Since $\dot{E}=0$ on the whole line $v=0$, we cannot claim ``strict decrease''; but
at any point of that line with $q\neq 0$ we have $\dot{v} = -(k/m)q \neq 0$: the
trajectory leaves immediately. Only the origin remains invariant there, and
\textbf{LaSalle's} invariance argument gives convergence.

% === FIGURE 13: spring without and with drag: (q,v) portrait and E(t) ===
\begin{figure}[H]
\centering
\includegraphics[width=0.95\textwidth]{figs/fig13_spring.pdf}
\caption{Left: without drag, closed orbits and constant $E$.
Right: with drag, energy falls and every trajectory converges.}
\label{fig:spring}
\end{figure}

\subsection{From energy to Lyapunov}

A differentiable scalar function $V$ is a Lyapunov candidate around a fixed point if
it is zero there and positive nearby. Compute $\dot{V} = \nabla V \cdot F$ using the
\emph{full} vector field. Strict negativity gives asymptotic stability (under the
stated domain conditions); if only $\dot{V}\leq 0$, one must identify the largest
invariant subset of $\{\dot{V}=0\}$. Energy is a statement about \emph{every}
trajectory, not about one simulated path.

\newpage

% ============================================================
\section{Chemistry forbids most vector fields}
\label{sec:chemistry}
% ============================================================

The spring was borrowed physics: its signed rotation is \emph{not} a closed
concentration model. An elementary mass-action network,

\begin{equation}
\sum_j \alpha_{ij} X_j \longrightarrow \sum_j \beta_{ij} X_j,
\qquad
v_i(x) = k_i \prod_j x_j^{\alpha_{ij}},
\label{eq:massaction}
\end{equation}

inherits restrictions that a generic vector field does not have:

\paragraph{Positivity.} If reaction $i$ consumes $X_j$, its flux contains the factor
$x_j$ and vanishes on the boundary $x_j=0$: consumption cannot push a concentration
negative. Every face of the nonnegative orthant points inward or is tangent. Pure
rotation $\dot{x}_1 = -x_2$, $\dot{x}_2 = x_1$ \emph{fails} as concentration
dynamics: on the face $x_1=0$ with $x_2>0$, $F_1 = -x_2 < 0$ and the field exits
the admissible region.

% === FIGURE 14: pure rotation violating positivity vs a field that satisfies it ===
\begin{figure}[H]
\centering
\includegraphics[width=0.8\textwidth]{figs/fig14_positivity.pdf}
\caption{Positivity excludes pure rotation as concentration dynamics: the
nonnegative orthant must be forward invariant.}
\label{fig:positivity}
\end{figure}

\paragraph{Independent parameters.} Distinct elementary reactions have distinct rate
constants. A behavior that requires two unrelated constants to match \emph{exactly}
is tuned, not structural.

\paragraph{Constrained kinetic form.} Mass action produces polynomials with signs
tied to what each reaction consumes and produces. Effective non-polynomial laws
(Hill, Michaelis--Menten) arise by \emph{eliminating fast states}: timescale
separation is a genuine route past the polynomial wall (the topic of lecture 6).

\paragraph{Conservation.} For $\dot{x} = Sv(x)$, every row vector $\ell$ with
$\ell^{\mathsf{T}} S = 0$ gives a conserved quantity $\ell^{\mathsf{T}} x =$
constant. Stability must then be judged \emph{inside} the corresponding
stoichiometric compatibility class: a zero eigenvalue normal to that class may be
reporting a conserved total, not a physical instability.

\newpage

% ============================================================
\section{Closing: the price list}
% ============================================================

Analysis did not replace simulation: it identified which questions can be settled
\emph{before} exact parameters are known, and which still require the bottom-level
model.

\begin{center}
\begin{tabular}{p{4.2cm} p{5.2cm} p{4.6cm}}
\hline
\textbf{function sought} & \textbf{structural price} & \textbf{what still has to be checked} \\
\hline
one persistent expression level & one restoring addition--removal crossing & domain, boundedness, recovery time \\
binary memory & three crossings; sufficiently steep cooperative response & basin sizes, noise-driven switching, reduction validity \\
commitment with hysteresis & two folds in a parameterized response & critical slowing, finite-copy noise \\
local recovery & a Hurwitz Jacobian on the relevant slice & whether the perturbation stays local \\
a biochemical clock & $\geq 2$ effective dimensions, rotational feedback, an isolated attracting orbit & boundedness, period, amplitude, robustness \\
global convergence & a Lyapunov function, monotonicity, contraction & every hypothesis on the invariant domain \\
\hline
\end{tabular}
\end{center}

The ``cell as a chemical plant'' question now has a more precise answer: positivity
keeps concentrations in their admissible region; dilution under constant growth
contributes self-limiting diagonal terms; saturation caps many production fluxes;
conserved pools confine motion to lower-dimensional slices. None of these alone is a
universal stability theorem ---cells deliberately switch and oscillate--- but together
they explain why unbounded behavior is not the default, and why interesting behavior
must be \emph{built}.

\begin{quote}
\centering
\emph{Simulation tells you what one cell did once. Structure tells you what every
cell of that design can and cannot do --- and that is the only kind of answer an
experiment can be designed against.}
\end{quote}

% ============================================================
% References
% ============================================================
\begin{thebibliography}{9}

\bibitem{scheffer2009}
M.~Scheffer \textit{et al.}, ``Early-warning signals for critical transitions,''
\textit{Nature} \textbf{461}, 53--59 (2009).

\bibitem{gardner2000}
T.~S. Gardner, C.~R. Cantor and J.~J. Collins, ``Construction of a genetic toggle
switch in \textit{Escherichia coli},'' \textit{Nature} \textbf{403}, 339--342 (2000).

\bibitem{elowitz2000}
M.~B. Elowitz and S.~Leibler, ``A synthetic oscillatory network of transcriptional
regulators,'' \textit{Nature} \textbf{403}, 335--338 (2000).

\bibitem{potvin2016}
L.~Potvin-Trottier \textit{et al.}, ``Synchronous long-term oscillations in a
synthetic gene circuit,'' \textit{Nature} \textbf{538}, 514--517 (2016).

\bibitem{hirsch}
M.~W. Hirsch, S.~Smale and R.~L. Devaney, \textit{Differential Equations, Dynamical
Systems, and an Introduction to Chaos}, 3rd ed., Academic Press (2013).

\bibitem{core}
F.~Xiao, \textit{CCBS 2026, Lecture 4 core page},
\url{https://chemaoxfz.github.io/assets/ccbs/2026fall/lecture04/core/}.

\end{thebibliography}

\end{document}
